\documentclass[11pt]{article}
\usepackage[T1]{fontenc}
\usepackage{lmodern}
\usepackage[margin=1in]{geometry}
\usepackage{amsmath,amssymb,amsthm,mathtools}
\usepackage{booktabs,array,enumitem,microtype}
\usepackage{xcolor}
\usepackage[colorlinks=true,linkcolor=blue!45!black,citecolor=blue!45!black,
  urlcolor=blue!45!black]{hyperref}
\hypersetup{pdftitle={Counting matroids on ten elements by aggregating extensions},
  pdfauthor={ainta}}
\setlength{\parindent}{1.2em}
\setlength{\parskip}{2pt}
\setlength{\emergencystretch}{1.5em}
\setlist{itemsep=2pt,topsep=4pt}
\newtheorem{theorem}{Theorem}[section]
\newtheorem{proposition}[theorem]{Proposition}
\newtheorem{lemma}[theorem]{Lemma}
\newtheorem{corollary}[theorem]{Corollary}
\theoremstyle{definition}
\newtheorem{example}[theorem]{Example}
\newcounter{algorithm}
\newenvironment{algorithm}[1]{%
  \refstepcounter{algorithm}\par\medskip\noindent
  \begin{minipage}{\linewidth}\hrule\medskip
  \textbf{Algorithm \thealgorithm. #1}\par\smallskip
  \begin{enumerate}[label=\arabic*.,leftmargin=2.2em,itemsep=2pt,topsep=0pt]
}{\end{enumerate}\medskip\hrule\end{minipage}\par\medskip}
\DeclareMathOperator{\Aut}{Aut}
\DeclareMathOperator{\cl}{cl}
\DeclareMathOperator{\type}{type}
\newcommand{\SP}{\mathrm{sp}}
\newcommand{\nsp}{\mathrm{nsp}}
\newcommand{\mc}{\operatorname{\#MC}}
\newcommand{\IS}{\mathsf{IS}}
\newcommand{\Cut}{\mathsf{Cut}}
\newcommand{\Part}{\mathcal P}
\newcommand{\Fl}{\mathcal F}
\newcommand{\ind}{\mathbf 1}
\newcommand{\uFive}{3\,222\,227\,959\,444}
\newcommand{\uTen}{3\,232\,000\,741\,644}
\newcommand{\uSparse}{2\,630\,337\,889\,305}
\newcommand{\uNonSparse}{591\,890\,070\,139}

\title{Counting matroids on ten elements\\by aggregating extensions}
\author{ainta}
\date{}

\begin{document}
\maketitle

\begin{abstract}
We give an exact algorithm for counting matroid isomorphism classes using
catalogues of smaller matroids. For each matroid and an automorphism, the
algorithm counts invariant modular cuts by branching on flats below the
hyperplanes and counting independent sets on the remaining hyperplanes.
Weighted sums over the smaller matroids recover the Burnside terms for
permutations with a fixed point; formulas expressing basis exchange supply
the remaining terms. We count sparse paving matroids separately using
Johnson graphs. We give pseudocode and correctness proofs for these
reductions and their composition. The implementation returns $\uFive$
classes of rank five on ten elements and
$\uTen$ classes across all ranks.
\end{abstract}

\section{The problem and the algorithm}

Let $u_{n,r}$ be the number of matroids of rank $r$ on $n$ elements up to
isomorphism, and let $L_{n,r}$ be the number on the labeled set
$[n]=\{1,\ldots,n\}$.
All matroids are included: loops, parallel elements, coloops, and disconnected
matroids are allowed. Isomorphism means a permutation of the ground set.
Duality gives $u_{n,r}=u_{n,n-r}$, but distinct duals remain distinct classes.

An extension $M$ by one element $e$ has a \emph{parent} $N=M\setminus e$.
The algorithm generates these smaller matroids and counts their extensions
in aggregate. Its two main reductions are
\[
\begin{aligned}
\text{one parent and one automorphism}
&\longrightarrow \text{branching on flats}
\longrightarrow \text{independent sets},\\
\text{counts of extensions}
&\longrightarrow \text{sums over parents}
\longrightarrow \text{isomorphism classes}.
\end{aligned}
\]
The first reduction avoids constructing the final matroids. The second
avoids a separate isomorphism test for each possible extension.

The main computation is rank five on ten elements. Sparse paving matroids
make up most of the classes in the resulting count, and their description
by independent sets of a Johnson graph gives a more efficient procedure
than the general extension constraints. We therefore count them separately
and restrict the general procedure to the remaining matroids:
\begin{equation}\label{eq:partition-result}
u_{10,5}=u^{\SP}_{10,5}+u^{\nsp}_{10,5}
       =\uSparse+\uNonSparse.
\end{equation}
The full algorithm also computes ranks three and four, uses elementary
formulas for ranks zero through two, and obtains the remaining ranks by
duality. Published counts are used for comparison, not as inputs.

The mathematical inputs are sets $\Part_{q,r}$ containing one representative
of each isomorphism class of matroids of rank $r$ on $q$ elements, with their
automorphism groups. The implementation generates the catalogues from the empty matroid using the
established extension generator of~\cite{MMIB,IC}; the counting procedure computes
their automorphism groups from the rank functions. The generator produces
$|\Part_{9,5}|=190\,214$, agreeing with~\cite{MR}.
The correspondence between modular cuts and extensions is due to
Crapo~\cite{Crapo}; canonical construction paths are standard~\cite{McKay}.
The subject of this note is
the aggregate counting algorithm built on these ingredients.

\paragraph{Notation.}
We use standard matroid terminology as in~\cite{Oxley,MR}. For a matroid
$N$, write $r_N$, $\cl_N$, $\Fl(N)$, and $\Aut(N)$ for its rank function,
closure operator, set of flats, and automorphism group. Superscripts
$\SP$ and $\nsp$ restrict a count to sparse paving matroids and matroids
that are not sparse paving, respectively. The counting symbols used below
are local to this note and are summarized here; each auxiliary quantity is
defined in its section.
\begin{center}
\small
\begin{tabular}{@{}lp{0.69\linewidth}@{}}
\toprule
Symbol & Meaning\\
\midrule
$u_{n,r},\ L_{n,r}$ & Numbers of unlabeled and labeled matroids\\
$F_{n,r}(\lambda)$ & Number of labeled matroids fixed by a permutation
of cycle type $\lambda$\\
$T_{n,r}(\mu)$ & Weighted sum over parents for a cycle type on $n-1$ elements\\
$c(N,g)$ & Number of nonempty modular cuts of $N$ invariant under $g$\\
$i(G),\ i_G(S)$ & Numbers of independent sets of $G$ and $G[S]$,
including the empty set\\
$G_{N,g}$ & Graph of conflicts between hyperplane orbits of a parent\\
$Q_\sigma$ & Graph of conflicts between orbits of subsets in a Johnson graph\\
$E(T)$ & Number of independent sets of $J(9,4)$ containing $T$\\
\bottomrule
\end{tabular}
\end{center}

\section{Counting extensions of one parent}\label{sec:cuts}

Let $N$ be a matroid on $E$, of rank $r=r_N(E)$.
A flat is a set equal to its closure; a hyperplane is a flat
of rank $r-1$. Two flats form a modular pair when
\begin{equation}\label{eq:modular-pair}
r_N(F)+r_N(H)=r_N(F\cup H)+r_N(F\cap H).
\end{equation}
Let $e\notin E$ be the element to be added.

\subsection{The Boolean formulation}

A modular cut $\mathcal C\subseteq\Fl(N)$ is upward closed and closed under
intersection of modular pairs. Crapo's correspondence~\cite{Crapo}, in the
form of~\cite[Theorem 7.2.3]{Oxley}, identifies
it with a unique extension $M$ satisfying $M\setminus e=N$. In this
correspondence,
\begin{equation}\label{eq:cut-meaning}
F\in\mathcal C\quad\Longleftrightarrow\quad e\in\cl_M(F).
\end{equation}
In particular, for $X\subseteq E$,
\begin{equation}\label{eq:extension-rank}
r_M(X)=r_N(X),\qquad
r_M(X\cup\{e\})=r_N(X)+1-\ind_{\{\cl_N(X)\in\mathcal C\}}.
\end{equation}
The empty cut gives a coloop extension and raises the rank by one.
Every nonempty cut contains $E$ and gives a rank-preserving extension.

Introduce $x_F\in\{0,1\}$ for every flat $F$. The nonempty cuts are exactly
the satisfying assignments of
\begin{align}
x_E&=1,\label{eq:top}\\
\neg x_F\lor x_H&\qquad(F\subsetneq H),\label{eq:upward}\\
\neg x_F\lor\neg x_H\lor x_{F\cap H}
&\qquad(F,H\text{ a modular pair}).\label{eq:cut-clause}
\end{align}
Thus counting extensions means counting these assignments, with the parent
labels fixed.

For $g\in\Aut(N)$, require the cut to be $g$-invariant by identifying flat
variables in the same orbit under $\langle g\rangle$. Simplify repeated
literals and delete tautological clauses. Denote the resulting formula by
$\Phi_{N,g}$. Its orbit variables retain a well-defined flat rank.

\begin{lemma}\label{lem:invariant-cuts}
The models of $\Phi_{N,g}$ are in bijection with rank-preserving extensions
of $N$ fixed by the permutation that acts as $g$ on $E$ and fixes $e$.
\end{lemma}
\begin{proof}
An assignment to orbit variables is exactly a $g$-invariant assignment to
flats. Equations~\eqref{eq:top}--\eqref{eq:cut-clause} say precisely that its
selected flats form a nonempty modular cut. The extension correspondence
is equivariant: equation~\eqref{eq:extension-rank} is preserved by $g$
exactly when the cut is preserved. Distinct cuts give distinct extensions.
\end{proof}

\subsection{Reduction to independent sets}

In this section, call a flat \emph{lower} if its rank is at most $r-2$.
Branch on the variables for orbits of lower flats and apply unit propagation
after every assignment. Once these variables are assigned, only variables
for hyperplane orbits can remain undecided.

For fixed $(N,g)$, define a simple graph $G_{N,g}$ whose vertices are the
$\langle g\rangle$-orbits of hyperplanes of $N$. Two distinct vertices
$O,O'$ are adjacent if and only if some $H\in O$ and $K\in O'$ satisfy
$r_N(H\cap K)=r-2$. We call this the \emph{conflict graph}.
An orbit can also contain two distinct hyperplanes with this property;
the modular cut clauses handle such an internal conflict.

\begin{lemma}[Hyperplane reduction]\label{lem:hyperplanes}
Suppose all variables for lower flats are assigned and unit propagation
has reached a consistent fixed point. Let $S$ be the set of undecided
hyperplane orbits. The satisfying completions of this assignment are in
bijection with the independent sets of $G_{N,g}[S]$.
\end{lemma}
\begin{proof}
Distinct hyperplanes span $N$, so~\eqref{eq:modular-pair} becomes
$r_N(H\cap K)=r-2$. Their intersection is a lower flat. If its variable
is true, upward closure forces both hyperplanes true. Consequently, whenever
both hyperplane variables remain undecided, their clause reduces
to $\neg x_H\lor\neg x_K$.

Every other clause has at most one undecided variable: its other flats are
lower or equal to $E$. Such a clause is either satisfied or has already
been resolved by unit propagation. If two conflicting hyperplanes lie in the
same orbit, a true intersection forces the orbit true, while a false
intersection forces it false. Such an orbit cannot remain undecided.
Thus the residual clauses are exactly the edge constraints of $G_{N,g}[S]$.
Selecting an independent set specifies exactly one completion of the cut.
\end{proof}

The graph $G_{N,g}$ is built once for each $(N,g)$; branching changes only
the set $S$. A parent of rank five on nine elements has at
most $\binom94=126$ hyperplanes and at most
$\sum_{j=0}^3\binom9j=130$ lower flats: a rank-$j$ flat is the closure of
some basis of size $j$. This bounds the graph size in the implementation.

\subsection{Counting independent sets}

For a simple graph $G$, let $i(G)$ count its independent sets, including
the empty set, and write $i_G(S)=i(G[S])$ for $S\subseteq V(G)$.
Write $N_G(v)$ for the open neighborhood of $v$.
Algorithm~\ref{alg:is} uses a cache keyed by the complete set $S$ for a
fixed graph. The cache may discard an entry whenever space is needed.

\begin{algorithm}{$\IS_G(S)$}\label{alg:is}
\item If $S=\varnothing$, return $1$. If $S$ is cached, return its value.
\item Save $S_0\gets S$. Remove all isolated vertices of $G[S]$;
      let $k$ be their number and $S$ the remaining set.
\item If $S=\varnothing$, set $a\gets 1$.
\item Otherwise, if component splitting finds a component $C\subsetneq S$,
      set $a\gets\IS_G(C)\,\IS_G(S\setminus C)$.
\item Otherwise, if $G[S]$ is a path on $t$ vertices, set $a\gets f_{t+2}$;
      if it is a cycle, set $a\gets f_{t-1}+f_{t+1}$.
\item Otherwise choose $v\in S$ of maximum degree and set
      $a\gets\IS_G(S\setminus\{v\})+
      \IS_G(S\setminus(\{v\}\cup N_G(v)))$.
\item Store $2^k a$ under the key $S_0$, discarding a cache entry if needed,
      and return $2^k a$.
\end{algorithm}

Here $f_0=0,f_1=1$ and $f_{t+2}=f_{t+1}+f_t$ define the Fibonacci numbers.
The rule for deciding when to search for components is an implementation
choice described in Section~\ref{sec:implementation}.

\begin{proposition}\label{prop:is}
Algorithm~\ref{alg:is} returns $i_G(S)$.
\end{proposition}
\begin{proof}
An isolated vertex can be independently included or excluded, giving $2^k$.
Choices on distinct components are independent, so their counts multiply.
Partitioning by whether $v$ is selected gives
\[
i_G(S)=i_G(S\setminus\{v\})+
i_G(S\setminus(\{v\}\cup N_G(v))).
\]
For paths, the same recurrence at an endpoint gives $f_{t+2}$; conditioning
on one vertex of a cycle gives $f_{t+1}+f_{t-1}$. Every recursive argument
is smaller, proving termination and correctness by induction on $|S|$.
Cache entries store previously established values of the same function.
Discarding an entry changes only the amount of repeated work.
\end{proof}

\subsection{Counting modular cuts}

We allow a set $\mathcal B$ of \emph{marked} orbit variables and a flag $b$ requiring
at least one of them to be true. Ordinary extension counting uses $b=0$.
Section~\ref{sec:nonsp} specifies the marks used to exclude sparse paving
extensions. Build $\Phi_{N,g}$ and $G=G_{N,g}$ once, and start with $x_E=1$.

\begin{algorithm}{$\Cut(\alpha,b)$ for fixed $\Phi_{N,g}$, $G=G_{N,g}$, and $\mathcal B$}
\label{alg:cut}
\item Propagate unit clauses from the partial assignment $\alpha$.
      If a contradiction occurs, return $0$.
\item If a marked variable is true, set $b\gets0$.
\item If an undecided lower variable $v$ exists, return
      $\Cut(\alpha\cup\{v=1\},b)+\Cut(\alpha\cup\{v=0\},b)$,
      using a separate state for each branch.
\item Let $S$ be the set of undecided hyperplane variables.
\item If $b=0$, return $\IS_G(S)$.
\item Set $a\gets0$. List the marked vertices in $S$ as $v_1,\ldots,v_t$.
\item For $j=1,\ldots,t$:
      set $S\gets S\setminus\{v_j\}$ and
      $a\gets a+\IS_G(S\setminus N_G(v_j))$.
\item Return $a$.
\end{algorithm}

The cache for $\IS_G$ is fresh for each $(N,g)$, since a key consisting only
of $S$ has meaning only after the graph $G$ is fixed.

\begin{proposition}[Correctness of the extension count]\label{prop:cuts}
Algorithm~\ref{alg:cut} counts exactly the models extending $\alpha$, subject
to the requirement that some marked variable is true when $b=1$.
\end{proposition}
\begin{proof}
Unit propagation forces values shared by every satisfying completion.
A contradiction gives no completions. Branching on $v$ partitions all
completions into two disjoint sets and decreases the number of undecided
lower variables. At a leaf, Lemma~\ref{lem:hyperplanes} and
Proposition~\ref{prop:is} give the unrestricted count.

If a marked variable must still be selected, partition the independent sets
by their first selected vertex in $v_1,\ldots,v_t$. For the $j$th part,
$v_1,\ldots,v_{j-1}$ are excluded and $v_j$ is included. The remaining
choices are precisely the independent sets of
$G[S_0\setminus(\{v_1,\ldots,v_j\}\cup N_G(v_j))]$, where $S_0$ was the
initial leaf set. These are the terms in lines 6--7. They are disjoint
and exhaust the required models; when $t=0$ their sum is zero.
\end{proof}

\begin{example}[Five extensions of a three-point line]\label{ex:line}
Take $N=U_{2,3}$. Its flats are $\varnothing$, the three singletons, and $E$.
If $x_{\varnothing}=1$, every flat is selected: the new element is a loop.
If $x_{\varnothing}=0$ and $g$ is the identity, the conflict graph is $K_3$, with four independent
sets. The empty set gives $U_{2,4}$, and each singleton makes $e$ parallel
to one existing point. Thus there are $1+4=5$ rank-preserving labeled
extensions. For the automorphism $(12)$, the orbit of the first two
hyperplanes has an internal conflict and is excluded when $x_{\varnothing}=0$.
Only the empty choice and the third hyperplane remain, giving $1+2=3$
invariant extensions. The coloop extension is separate and has rank three.
\end{example}

\section{From parent counts to isomorphism classes}\label{sec:symmetry}

Let $F_{n,r}(\lambda)$ be the number of labeled matroids of rank $r$ fixed by
one permutation of cycle type $\lambda\vdash n$. If $a_j(\lambda)$ is the
number of parts equal to $j$, set
\[
z_\lambda=\prod_{j\ge1}j^{a_j(\lambda)}a_j(\lambda)!.
\]
There are $n!/z_\lambda$ permutations of this type. Burnside's lemma gives
\begin{equation}\label{eq:burnside}
u_{n,r}=\sum_{\lambda\vdash n}\frac{F_{n,r}(\lambda)}{z_\lambda}.
\end{equation}
The identity term $F_{n,r}(1^n)$ is $L_{n,r}$.

\subsection{Permutations with a fixed element}

Put $q=n-1$. For $N\in\Part_{q,r}$ and $g\in\Aut(N)$, let $c(N,g)$ be
the number of nonempty modular cuts invariant under $g$, computed by
Algorithm~\ref{alg:cut} with $b=0$. For every $\mu\vdash q$, define the
following weighted sum over parents:
\begin{equation}\label{eq:parent-sum}
\begin{split}
T_{n,r}(\mu)={}&
\sum_{N\in\Part_{q,r}}\frac{q!}{|\Aut(N)|}
 \sum_{\substack{g\in\Aut(N)\\\type(g)=\mu}}c(N,g)\\
&+\sum_{N\in\Part_{q,r-1}}\frac{q!}{|\Aut(N)|}
 \#\{g\in\Aut(N):\type(g)=\mu\}.
\end{split}
\end{equation}
The second line counts coloop extensions. Write $\mu\cup(1)$ for adjoining
one part of size one.

\begin{theorem}[Counting by deletion]\label{thm:aggregation}
For every $\mu\vdash n-1$,
\begin{equation}\label{eq:fixed-from-sum}
F_{n,r}(\mu\cup(1))=\frac{z_\mu T_{n,r}(\mu)}{(n-1)!}.
\end{equation}
\end{theorem}
\begin{proof}
Fix the added label $e=n$. Count pairs $(M,g)$, where $M$ is a labeled
matroid of rank $r$ on $[n]$, $g$ is a permutation of $[q]$ of type $\mu$,
and $g$ extended by fixing $e$ preserves $M$.

Delete $e$. If $e$ is not a coloop, the parent has rank $r$, is preserved
by $g$, and Lemma~\ref{lem:invariant-cuts} supplies $c(N,g)$ choices.
If $e$ is a coloop, the parent has rank $r-1$ and there is exactly one
extension. Each parent type has $q!/|\Aut(N)|$ labelings. Summing by parent
therefore gives precisely~\eqref{eq:parent-sum}.

Alternatively, there are $q!/z_\mu$ choices of $g$. They are conjugate in
$S_q$, so each fixes $F_{n,r}(\mu\cup(1))$ matroids. Equating the two
counts proves~\eqref{eq:fixed-from-sum}.
\end{proof}

The contribution to~\eqref{eq:burnside} from permutations with a fixed point is
\begin{equation}\label{eq:fixed-contribution}
\sum_{\mu\vdash q}\frac{T_{n,r}(\mu)}{q!\,(a_1(\mu)+1)},
\end{equation}
because $z_{\mu\cup(1)}=(a_1(\mu)+1)z_\mu$. This factor is the precise
correction for adding a distinguished fixed element. Summing extension
orbits directly would instead count pointed matroids $(M,e)$, with one
occurrence per element orbit of $M$.

\paragraph{Conjugacy classes.}
If $g'=hgh^{-1}$ for $h\in\Aut(N)$, the map
$\mathcal C\mapsto h(\mathcal C)$ bijects the invariant cuts for $g$ and
$g'$. Thus each inner sum in~\eqref{eq:parent-sum} may be evaluated
using one representative per conjugacy class of $\Aut(N)$, multiplied by
its class size. Equal cycle types in $S_q$ alone do not establish this
bijection for a fixed parent: the conjugating permutation must preserve $N$.

\subsection{Permutations without a fixed element}

For a representative $\sigma$ of type $\lambda\vdash n$, introduce one
Boolean variable $y_O$ for every orbit $O$ of $r$-subsets under
$\langle\sigma\rangle$. Let $[B]$ denote the orbit of $B$. Use the formula
\begin{align}
\Psi_{n,r,\sigma}={}&\left(\bigvee_O y_O\right)\ \land
\bigwedge_{\substack{B,C\in\binom{[n]}r\\a\in B\setminus C}}
\left(\neg y_{[B]}\lor\neg y_{[C]}\lor
 \bigvee_{b\in C\setminus B}y_{[(B\setminus\{a\})\cup\{b\}]}\right).
\label{eq:bases}
\end{align}
Write $\mc(\Psi)$ for its exact number of satisfying assignments.

\begin{proposition}\label{prop:bases}
$\mc(\Psi_{n,r,\sigma})=F_{n,r}(\lambda)$.
\end{proposition}
\begin{proof}
An assignment selects a unique $\sigma$-invariant family of $r$-subsets.
The first clause makes it nonempty. The remaining clauses are exactly the
basis exchange axiom, which characterizes the bases of matroids of rank
$r$~\cite[Section 1.2]{Oxley}. Conversely, the bases of every fixed matroid give one satisfying
assignment. No auxiliary variables occur, so the correspondence preserves
the number of models without any projection factor.
\end{proof}

For a hard formula we choose $k$ variables and count all $2^k$ restrictions.
Every model has exactly one assignment to these variables, hence
\begin{equation}\label{eq:cubes}
\mc(\Psi)=\sum_{\eta\in\{0,1\}^k}\mc(\Psi\land\eta).
\end{equation}
These jobs can be scheduled independently. At $n=10$ there are $42$ cycle
types: $30$ with a fixed point and $12$ without one. The latter are the
only types for which we count models of $\Psi_{n,r,\sigma}$.

\begin{algorithm}{Full rank count $\mathsf{RankCount}(n,r)$, $1\le r<n$}
\label{alg:rank}
\item Set $q\gets n-1$ and $T_{n,r}(\mu)\gets0$ for every $\mu\vdash q$.
\item For each $N\in\Part_{q,r}$, put $w\gets q!/|\Aut(N)|$.
      For every representative $g$ of a conjugacy class in $\Aut(N)$ with class
      size $s$, compute $c\gets\Cut(x_E=1,0)$ and add $ws c$ to
      $T_{n,r}(\type(g))$.
\item For each $N\in\Part_{q,r-1}$ and each such class $(s,g)$,
      add $q!s/|\Aut(N)|$ to $T_{n,r}(\type(g))$.
\item For each $\mu\vdash q$, set
      $F_{n,r}(\mu\cup(1))\gets z_\mu T_{n,r}(\mu)/q!$, checking exact division.
\item For each $\lambda\vdash n$ with no part $1$, choose $\sigma$ of that
      type and set $F_{n,r}(\lambda)\gets\mc(\Psi_{n,r,\sigma})$.
\item Set $A\gets\sum_{\lambda\vdash n}(n!/z_\lambda)F_{n,r}(\lambda)$.
      Return $(A/n!,F_{n,r}(1^n))$, checking exact division.
\end{algorithm}

\begin{theorem}\label{thm:rank}
Given complete parent catalogues and their automorphism groups,
Algorithm~\ref{alg:rank}, with exact formula counts, returns $(u_{n,r},L_{n,r})$.
\end{theorem}
\begin{proof}
Proposition~\ref{prop:cuts} supplies the correct invariant cut counts.
Summing over conjugacy classes gives the inner sums of
\eqref{eq:parent-sum}. Theorem~\ref{thm:aggregation} supplies every
term for a permutation with a fixed point, and Proposition~\ref{prop:bases} supplies every remaining
term. These two sets of cycle types partition all of $S_n$.
Equation~\eqref{eq:burnside} proves the unlabeled count; the identity fixes
every labeled matroid and gives the second output.
\end{proof}

\section{Excluding sparse paving extensions}\label{sec:nonsp}

A matroid of rank $r$ is paving if every circuit has size at least $r$, and
sparse paving if both it and its dual are paving. For rank five on ten
elements this means that every circuit and every cocircuit has size
at least five. Equivalently, it is paving and every hyperplane has at most
five elements, since the complement of a hyperplane is a cocircuit.

In this section $N=M\setminus e$ has rank five on nine elements, and $M$
also has rank five. If $N$ is not sparse paving, neither is $M$: a circuit
of $N$ of size at most four is also a circuit of $M$, and a hyperplane of
$N$ of size at least six lies in a hyperplane of $M$ of size at least six.

\begin{lemma}[The marked flats]\label{lem:marked}
Suppose $N$ is sparse paving and $\mathcal C$ is the nonempty modular cut
determining $M$.
The extension $M$ is not sparse paving if and only if $\mathcal C$ contains
either a flat of rank at most three or a hyperplane of size five.
\end{lemma}
\begin{proof}
If $F\in\mathcal C$ has rank at most three, choose a basis $I$ of $F$.
Then $e\in\cl_M(I)$, so $I\cup\{e\}$ contains a circuit through $e$
of size at most four. If a hyperplane $H$ of size five belongs to
$\mathcal C$, its closure in $M$ is $H\cup\{e\}$, a hyperplane of size six.
Its complement is a cocircuit of size four. Either case
violates sparse paving.

Conversely, if $M$ has a circuit $D$ of size at most four, then $e\in D$
because $N$ is paving. The flat $\cl_N(D\setminus\{e\})$ has rank at
most three and belongs to $\mathcal C$. Otherwise $M$ is paving, so failure
of sparse paving gives a hyperplane $K$ with $|K|\ge6$. It must contain
$e$: a hyperplane of $M$ avoiding $e$ is a hyperplane of $N$, whose size
is at most five. The flat $H=K\setminus\{e\}$ of $N$ has at least five
elements. Since every subset of $N$ of size four is independent,
$r_N(H)=4$. Thus $H$ is a hyperplane, $|H|=5$, and
$r_M(H\cup\{e\})=r_N(H)$, proving $H\in\mathcal C$.
\end{proof}

Let $c^{\nsp}(N,g)$ be the number of nonempty modular cuts invariant under
$g$ whose extensions are not sparse paving. If $N$ is not sparse paving, use
Algorithm~\ref{alg:cut} with $b=0$. Otherwise mark the orbit variables
containing the flats in Lemma~\ref{lem:marked}, and use $b=1$. These marks
are invariant under $\Aut(N)$, so the reduction to conjugacy classes still
applies. Proposition~\ref{prop:cuts} proves that the omitted branch is
exactly the sparse paving family.

\subsection{The coloop contribution at the middle rank}

Every coloop extension of rank five on ten elements fails to be sparse
paving, because a coloop is a cocircuit of size one. Duality bijects
$\Part_{9,5}$ and $\Part_{9,4}$ and preserves their automorphism groups
as permutation groups. Consequently the two parent sums can be paired:
\begin{equation}\label{eq:nsp-sum}
T^{\nsp}_{10,5}(\mu)=
\sum_{N\in\Part_{9,5}}\frac{9!}{|\Aut(N)|}
 \sum_{\substack{g\in\Aut(N)\\\type(g)=\mu}}
       \bigl(c^{\nsp}(N,g)+1\bigr),\qquad \mu\vdash9.
\end{equation}
The $+1$ is the coloop extension of the dual parent. This pairing is used
only at $(n,r)=(10,5)$; the general rank calculation keeps the two parent
sets in~\eqref{eq:parent-sum} separate.

Applying Theorem~\ref{thm:aggregation} to the family of matroids that are
not sparse paving gives
\begin{equation}\label{eq:nsp-fixed}
F^{\nsp}_{10,5}(\lambda)=
\begin{cases}
z_\mu T^{\nsp}_{10,5}(\mu)/9!,&\lambda=\mu\cup(1),\\
\mc(\Psi_{10,5,\sigma})-F^{\SP}_{10,5}(\lambda),&1\notin\lambda.
\end{cases}
\end{equation}
Thus
\begin{equation}\label{eq:nsp-total}
u^{\nsp}_{10,5}=\sum_{\lambda\vdash10}
                  \frac{F^{\nsp}_{10,5}(\lambda)}{z_\lambda}.
\end{equation}
Here $\sigma$ has cycle type $\lambda$. For permutations with a fixed point,
the procedure omits all extensions in which every marked variable is false.
For the twelve derangement types, it counts all matroids using basis exchange
and subtracts the number that are sparse paving.

\section{The sparse paving calculation}\label{sec:sparse}

For $1\le r<n$, let $J(n,r)$ be the Johnson graph: its vertices are the
$r$-subsets of $[n]$, with two vertices adjacent if and only if their
intersection has size $r-1$. The circuit-hyperplanes of a sparse paving
matroid of rank $r$ form an independent set of this graph, and every
independent set specifies such a matroid~\cite{MR,PvdP}.
These sets are also called \emph{stable sets} in the literature. Thus
$L^{\SP}_{n,r}=i(J(n,r))$. This description removes the need to branch on
lower flats and is the reason for counting this family separately.
All orbits below are under permutations of the ground set. When $n=2r$,
complementation also preserves $J(n,r)$; we do not quotient by it, because
it corresponds to matroid duality.

\subsection{The identity permutation}

Distinguish $e=n$ and split the family of circuit-hyperplanes into $A$, the blocks
containing $e$ with $e$ removed, and $B$, the blocks avoiding $e$.
Define
\[
\Gamma^+(A)=\left\{b\in\binom{[n-1]}r:\exists a\in A,\ a\subset b\right\}.
\]
Then $A$ is independent in $J(n-1,r-1)$, $B$ is independent in $J(n-1,r)$,
and no block in one part is adjacent to a block in the other exactly when
$B\cap\Gamma^+(A)=\varnothing$.

\begin{proposition}[Splitting at one element]\label{prop:split}
If $\mathcal A$ contains one representative from every $S_{n-1}$-orbit
of independent sets of $J(n-1,r-1)$, then
\begin{equation}\label{eq:split}
L^{\SP}_{n,r}=\sum_{A\in\mathcal A}
       \frac{(n-1)!}{|\Aut(A)|}\,
       i\bigl(J(n-1,r)-\Gamma^+(A)\bigr).
\end{equation}
\end{proposition}
\begin{proof}
Every family of circuit-hyperplanes has a unique pair $(A,B)$. For a fixed
labeled $A$, its possible $B$ are the independent sets of the displayed
residual graph. Relabeling $A$ gives an isomorphic residual graph, so all
$(n-1)!/|\Aut(A)|$ members of its orbit have the same number of extensions.
Summing over the parent orbits counts every labeled family once.
\end{proof}

\subsection{Completion counts for small parents}

A parent with few blocks forbids few vertices, leaving a large residual
graph. These cases are expensive for direct counting. For $(n,r)=(10,5)$,
we instead use inclusion--exclusion and a table of completion counts when
the parent has at most five blocks.

Complementation on $[9]$ identifies $J(9,5)$ with $G=J(9,4)$. A parent
block $a\in\binom{[9]}4$ forbids its five supersets of size five.
After complementation these become precisely
\[
D_a=\binom{[9]\setminus a}{4},
\]
the five subsets of size four of $[9]\setminus a$. They form a clique in
$G$. Two distinct blocks of $A$ cannot share a superset of size five,
because $A$ is independent in $J(9,4)$. Thus the forbidden set
$D=\bigcup_{a\in A}D_a$ is a union of $|A|$ disjoint cliques.

For an independent set $T$ of $G$, define
$E(T)=\#\{H\subseteq V(G):H\text{ independent},\ T\subseteq H\}$.
Let $o(H)=9!/|\Aut(H)|$ and
$m_T(H)=\#\{T'\subseteq H:T'\cong T\}$, with isomorphism under $S_9$.

\begin{proposition}[Counting completions and applying inclusion--exclusion]\label{prop:completion}
Given a catalogue $\mathcal H$ containing one representative of every $S_9$-orbit of independent sets of
$G$, we have
\begin{align}
E(T)&=\frac{\sum_{H\in\mathcal H}o(H)m_T(H)}{o(T)},
\label{eq:completion}\\
i(G-D)&=\sum_{\substack{T\subseteq D\\T\text{ independent}}}
                  (-1)^{|T|}E(T).
\label{eq:avoid}
\end{align}
If $D$ is covered by $k$ cliques, all terms with $|T|>k$ vanish.
\end{proposition}
\begin{proof}
For~\eqref{eq:completion}, count pairs $(H,T')$ over all labeled independent
sets $H$, with $T'\subseteq H$ and $T'\cong T$. Summing first over $H$
gives the numerator. Each of the $o(T)$ copies of $T$ has the same number
$E(T)$ of completions, giving the denominator.

For~\eqref{eq:avoid}, a fixed independent set $H$ contributes
$\sum_{T\subseteq H\cap D}(-1)^{|T|}$, which equals one when
$H\cap D=\varnothing$ and zero otherwise. Every such $T$ is independent.
An independent set can meet each covering clique in at most one vertex,
so its size is at most $k$.
\end{proof}

The implementation tabulates $E(T)$ for $|T|\le5$ from the freshly generated
catalogue of sparse paving matroids of rank four on nine elements. This handles all $251$ parent types
with $|A|\le5$. Larger parents use Algorithm~\ref{alg:is} on the residual
graph. This is a choice between two exact formulas for the same count.

\subsection{The other permutations}

For $\sigma\in S_n$, an invariant family of circuit-hyperplanes is a union of
$\langle\sigma\rangle$-orbits of $r$-subsets. Delete any orbit containing
two members that are adjacent in $J(n,r)$. Form a graph $Q_\sigma$ on the
remaining orbits, with two vertices adjacent if and only if some member
of one orbit is adjacent in $J(n,r)$ to a member of the other.
Selecting an independent set of $Q_\sigma$ is equivalent to selecting a $\sigma$-invariant independent
set of $J(n,r)$; hence
\begin{equation}\label{eq:sp-fixed}
F^{\SP}_{n,r}(\type(\sigma))=i(Q_\sigma).
\end{equation}

\begin{algorithm}{Counting invariant sparse paving matroids at $(n,r)=(10,5)$}
\label{alg:sparse}
\item Generate the complete catalogue $\mathcal H$ of independent sets of
      $J(9,4)$ modulo $S_9$, including their automorphism orders.
\item Compute $E(T)$ for the types with $|T|\le5$
      using~\eqref{eq:completion}.
\item For each $A\in\mathcal H$, compute its residual count by
      \eqref{eq:avoid} if $|A|\le5$, and by $\IS$ otherwise.
\item Set $F^{\SP}_{10,5}(1^{10})$ to the sum~\eqref{eq:split}.
\item For every other $\lambda\vdash10$, choose a permutation $\sigma$ of
      type $\lambda$, construct $Q_\sigma$, and set
      $F^{\SP}_{10,5}(\lambda)\gets\IS_{Q_\sigma}(V(Q_\sigma))$.
\item Return the table $F^{\SP}_{10,5}$ and
      $u^{\SP}_{10,5}=\sum_{\lambda\vdash10}F^{\SP}_{10,5}(\lambda)/z_\lambda$.
\end{algorithm}

\begin{corollary}\label{cor:sparse}
Algorithm~\ref{alg:sparse} returns the numbers of fixed sparse paving
matroids and their unlabeled total.
\end{corollary}
\begin{proof}
Propositions~\ref{prop:split} and~\ref{prop:completion} give the identity
term. The bijection given by selecting orbits proves~\eqref{eq:sp-fixed} for every
other term. Proposition~\ref{prop:is} evaluates these graph counts, and
Burnside's lemma gives the final sum.
\end{proof}

\section{Counting all matroids on ten elements}

For the three smallest ranks, direct formulas are sufficient. For $n\ge2$,
\begin{align}
(u_{n,0},L_{n,0})&=(1,1),\nonumber\\
(u_{n,1},L_{n,1})&=(n,2^n-1),\label{eq:small-ranks}\\
u_{n,2}&=\sum_{s=2}^n\bigl(p(s)-1\bigr),\qquad
L_{n,2}=B_{n+1}-2^n.\nonumber
\end{align}
Here $p(s)$ is the integer partition number and $B_s$ is the Bell number.
Rank one is determined by its nonempty set of nonloops. A matroid of rank two
is determined by its loops and a partition of the nonloops into at least
two parallel classes. Up to isomorphism, the class sizes are an integer
partition of some $s\le n$ with at least two parts. For labeled counting,
adjoin a marker to the loop block and partition the resulting $n+1$
elements. The forbidden cases with zero or one nonloop block number
$1+(2^n-1)=2^n$. This proves~\eqref{eq:small-ranks}.

\begin{algorithm}{All matroids on ten elements}
\label{alg:ten}
\item Generate the parent catalogues $\Part_{9,r}$ needed below, with their
      automorphism groups, and the sparse paving catalogue for
      Algorithm~\ref{alg:sparse}.
\item Run Algorithm~\ref{alg:sparse} to obtain $F^{\SP}_{10,5}(\lambda)$ for all
      $\lambda\vdash10$.
\item Set $T^{\nsp}_{10,5}(\mu)\gets0$ for all $\mu\vdash9$.
      For each $N\in\Part_{9,5}$ and each automorphism conjugacy class
      $(s,g)$, compute $c^{\nsp}(N,g)$ as in Section~\ref{sec:nonsp} and add
      $9!s(c^{\nsp}(N,g)+1)/|\Aut(N)|$ to $T^{\nsp}_{10,5}(\type(g))$.
\item Obtain $F^{\nsp}_{10,5}(\lambda)$ from~\eqref{eq:nsp-fixed}, and set
      $F_{10,5}(\lambda)\gets F^{\SP}_{10,5}(\lambda)+F^{\nsp}_{10,5}(\lambda)$.
      Check the exact divisions and nonnegativity.
\item Set $u_{10,5}\gets\sum_{\lambda\vdash10}F_{10,5}(\lambda)/z_\lambda$
      by integer Burnside accumulation, and $L_{10,5}\gets F_{10,5}(1^{10})$.
\item Compute ranks $3$ and $4$ with Algorithm~\ref{alg:rank} and ranks
      $0,1,2$ with~\eqref{eq:small-ranks}.
\item Return
      $\bigl(u_{10,5}+2\sum_{r=0}^4u_{10,r},\;
      L_{10,5}+2\sum_{r=0}^4L_{10,r}\bigr)$.
\end{algorithm}

\begin{theorem}[Correctness of the combined count]\label{thm:ten}
With complete parent catalogues, their exact automorphism groups, and exact
counts of the formulas expressing basis exchange, Algorithm~\ref{alg:ten} returns the numbers
of unlabeled and labeled matroids on ten elements, across all ranks.
\end{theorem}
\begin{proof}
Corollary~\ref{cor:sparse} gives the numbers of fixed sparse paving matroids.
Lemma~\ref{lem:marked} and Proposition~\ref{prop:cuts} count the invariant
extensions that are not sparse paving for each parent of rank five.
Equation~\eqref{eq:nsp-sum} also includes every coloop extension through
the dual parent of rank four. Theorem~\ref{thm:aggregation}, applied to this
family, therefore gives its fixed counts for every permutation with a
fixed point. Proposition~\ref{prop:bases} gives the full derangement
terms, from which the sparse paving terms are subtracted.

The two families are disjoint and exhaustive, so their fixed counts add to
$F_{10,5}(\lambda)$ for every cycle type. Burnside's lemma and the identity
term give the two counts at rank five. Theorem~\ref{thm:rank} and
\eqref{eq:small-ranks} establish ranks zero through four. Finally, duality
is a bijection on both labeled matroids and isomorphism classes, pairing
each rank below five with a distinct rank above five. The last line sums
these disjoint ranks exactly once.
\end{proof}

\section{Implementation and computation}

\subsection{Implementation}\label{sec:implementation}

To compute $\Aut(N)$, the implementation backtracks over bijections of the parent
ground set. A candidate image must have the same counts of containing
subsets of each size and rank. After each assignment, it checks rank
preservation on every subset of the assigned domain. Every automorphism
passes both tests, while every completed bijection preserves the full rank
function; thus the returned group is exact. Its conjugacy classes are
formed by explicitly collecting $hgh^{-1}$ for $h\in\Aut(N)$.

The C++ implementation of Algorithm~\ref{alg:cut} chooses an undecided
lower variable with the most clause occurrences. It records assignments
on a trail and undoes them after each branch. Algorithm~\ref{alg:is}
stores vertex sets as bit masks and attempts component decomposition when
fewer than $45$ vertices remain or the maximum degree is below $5$.
Otherwise it branches directly.

Persistent processes take parents from a shared queue. For each parent,
a process evaluates one representative of every automorphism conjugacy
class, with a fresh graph cache for each representative. It returns the
parent's contributions to the sums $T_{n,r}(\mu)$ or
$T^{\nsp}_{10,5}(\mu)$. These sums and the final reductions use
arbitrary-precision integers. The code calls these arrays \texttt{moments};
mathematically they are the weighted sums defined above.

The formulas for derangements use one variable per orbit of subsets and no auxiliary
variables. Ganak~\cite{Ganak} is run in exact mode. For the type $2^5$ at
ranks four and five, the driver fixes eight variables of highest clause
occurrence and sums all $256$ resulting cubes, using~\eqref{eq:cubes}.
Completed results are saved so that an interrupted computation can resume.
A job that times out contributes no count until it completes.

The source correspondence is:
\begin{center}
\small
\begin{tabular}{@{}ll@{}}
\toprule
Part of the algorithm & Source file in the repository\\
\midrule
Counting independent sets & \texttt{src/split\_count.cpp}\\
Invariant modular cuts & \texttt{src/extension\_worker.cpp}\\
Completion counts & \texttt{src/small\_extensions.hpp}\\
Counting sparse paving matroids & \texttt{scripts/count\_sparse.py}\\
Basis exchange formulas & \texttt{scripts/export\_symmetry.py}\\
Burnside's lemma and sum over ranks & \texttt{scripts/count.py}\\
\bottomrule
\end{tabular}
\end{center}

\subsection{Results}

Table~\ref{tab:ranks} gives the outputs by rank. The full sums are
\begin{equation}\label{eq:totals}
\boxed{u_{10}=\uTen},\qquad
\boxed{L_{10}=11\,727\,995\,799\,397\,397\,461}.
\end{equation}
The unlabeled counts at ranks zero through four agree with
\cite[Table 1]{MMIB}. The value at rank five is a result of the present
computation; the full count at this rank has not yet been replicated
independently.

\begin{table}[htbp]
\centering
\caption{Computed counts on ten elements. Ranks $10-r$ have the same values.}
\label{tab:ranks}
\begin{tabular}{@{}rrr@{}}
\toprule
Rank $r$ & Unlabeled $u_{10,r}$ & Labeled $L_{10,r}$\\
\midrule
0 & 1 & 1\\
1 & 10 & 1\,023\\
2 & 128 & 677\,546\\
3 & 10\,037 & 16\,869\,404\,175\\
4 & 4\,886\,380\,924 & 17\,716\,560\,870\,661\,325\\
5 & \uFive & 11\,692\,562\,643\,915\,909\,321\\
\bottomrule
\end{tabular}
\end{table}

The calculation at rank five separates as follows. For each column, the
middle row is the contribution of all nonidentity permutations to the
Burnside numerator:
$\sum_{\lambda\ne1^{10}}(10!/z_\lambda)F^{\SP}_{10,5}(\lambda)$ or
the same sum with superscript $\nsp$. Adding the first two rows and
dividing by $10!$ gives the last row.
\begin{center}
\small
\begin{tabular}{@{}lrr@{}}
\toprule
 & Sparse paving & Not sparse paving\\
\midrule
Identity
 & 9\,544\,867\,699\,823\,218\,639
 & 2\,147\,694\,944\,092\,690\,682\\
Nonidentity permutations
 & 102\,432\,886\,765\,361 & 155\,742\,427\,712\,518\\
Unlabeled count & \uSparse & \uNonSparse\\
\bottomrule
\end{tabular}
\end{center}
Mayhew--Royle~\cite[Section 9.2]{MR} estimated approximately
$2.65\times10^{12}$ sparse paving classes of rank five on ten elements.
Our sparse paving count is close to that estimate, which is not an exact
value against which the computation can be verified.

The complete tables of numbers of fixed matroids are recorded in
\texttt{results/fixed\_terms.json}. Appendix~\ref{app:derangements} lists
the twelve counts for derangements at rank five.

The recorded full run took $389.16$ seconds on $64$ physical cores of two
AMD EPYC 9354 processors. This includes generating the parents from the
empty matroid ($3.83$ seconds), the sparse paving calculation ($124.85$
seconds), ranks three and four ($6.92$ and $76.72$ seconds), and the
remaining work at rank five ($176.47$ seconds). Compilation and the separate
validation runs are excluded. The C++ programs were compiled with GCC
13.3.0 and \texttt{-O3 -march=native -std=c++17}.

\subsection{Checks and reproduction}

The generated catalogues reproduce all $55$ entries indexed by size and rank through
nine elements in~\cite{MR}, including $383\,172$ classes on nine elements.
Across sizes zero through nine, all $385\,370$ canonical rank arrays match
the reference catalogue~\cite{Catalogue}. The rank axioms, absence of
isomorphic duplicates, and closure under duality were checked directly.
For the $190\,214$ parents of rank five on nine elements, the orders of the
automorphism groups were also checked by a separate program. Relabeling
each parent preserves its contribution to every sum $T^{\nsp}_{10,5}(\mu)$.

The aggregate algorithm itself reproduces the known values
\[
\begin{array}{c|rrrrrr}
(n,r)&(6,3)&(7,3)&(8,3)&(8,4)&(9,3)&(9,4)\\ \hline
u_{n,r}&38&108&325&940&1\,275&190\,214.
\end{array}
\]
These cases exercise both the sums over parents and the formulas for
derangements. A fresh computation on ten elements with regenerated parents
reproduces all $126$ numbers of fixed matroids at ranks three, four, and
five. The result and check
summaries are in \texttt{results/counts.json},
\texttt{results/validation.json}, and \texttt{results/provenance.json}.

\noindent\begin{minipage}{\linewidth}
The implementation is in the \texttt{matroid-count} repository~\cite{Code},
revision \texttt{d1aaa1468a01}. From its root, on Linux with Python 3.10+,
GCC, Make, and Boost headers, run:
\begin{verbatim}
python3 scripts/setup.py --audit
make -j4
make audit
make test
python3 scripts/count.py --workdir runs/n10 --jobs 64
python3 scripts/validate.py --parents runs/n10/parents \
    --ten-run runs/n10 --out runs/n10/validation --jobs 8
\end{verbatim}
\end{minipage}
\par\smallskip
The computed total is written to \texttt{runs/n10/total.json}.
The generated parent ledgers, formulas, and solver logs remain in that
run directory. The source tree contains the compact result summaries.

\appendix
\section{Derangement terms at rank five}\label{app:derangements}

Each entry in Table~\ref{tab:derangements} is the number of labeled matroids
fixed by one permutation of the indicated type. The last column is obtained
by subtraction. Its contribution to the unlabeled count is the sum of that
column divided, row by row, by $z_\lambda$. Types are ordered by increasing
number of parts, then in decreasing lexicographic order, with the parts of
each partition written in decreasing order.

\begin{table}[htbp]
\centering
\caption{All twelve derangement types in $S_{10}$ at rank five.}
\label{tab:derangements}
\begin{tabular}{@{}lrrr@{}}
\toprule
Cycle type & All matroids & Sparse paving & Not sparse paving\\
\midrule
$10$      & 78 & 67 & 11\\
$8,2$     & 637 & 389 & 248\\
$7,3$     & 28 & 1 & 27\\
$6,4$     & 746 & 132 & 614\\
$5^2$     & 24\,566 & 21\,969 & 2\,597\\
$6,2^2$   & 5\,620 & 1\,146 & 4\,474\\
$5,3,2$   & 129 & 4 & 125\\
$4^2,2$   & 150\,791 & 53\,705 & 97\,086\\
$4,3^2$   & 1\,022 & 16 & 1\,006\\
$4,2^3$   & 792\,347 & 34\,257 & 758\,090\\
$3^2,2^2$ & 104\,620 & 16\,542 & 88\,078\\
$2^5$     & 584\,746\,603 & 217\,942\,377 & 366\,804\,226\\
\bottomrule
\end{tabular}
\end{table}

\begin{thebibliography}{10}
\small
\bibitem{Crapo}
H.~H. Crapo,
Single-element extensions of matroids,
\emph{J. Res. Nat. Bur. Standards Sect. B} \textbf{69B} (1965), 55--65.
\href{https://doi.org/10.6028/jres.069B.003}{doi:10.6028/jres.069B.003}.

\bibitem{MR}
D. Mayhew and G.~F. Royle,
Matroids with nine elements,
\emph{J. Combin. Theory Ser. B} \textbf{98} (2008), 415--431.
\href{https://doi.org/10.1016/j.jctb.2007.07.005}{doi:10.1016/j.jctb.2007.07.005}.

\bibitem{McKay}
B.~D. McKay,
Isomorph-free exhaustive generation,
\emph{J. Algorithms} \textbf{26} (1998), 306--324.
\url{https://users.cecs.anu.edu.au/~bdm/papers/orderly.pdf}.

\bibitem{MMIB}
Y. Matsumoto, S. Moriyama, H. Imai, and D. Bremner,
Matroid enumeration for incidence geometry,
\emph{Discrete Comput. Geom.} \textbf{47} (2012), 17--43.
\href{https://doi.org/10.1007/s00454-011-9388-y}{doi:10.1007/s00454-011-9388-y}.

\bibitem{Oxley}
J. Oxley,
\emph{Matroid Theory}, second edition,
Oxford University Press, 2011.

\bibitem{PvdP}
R. Pendavingh and J. van der Pol,
On the number of matroids compared to the number of sparse paving matroids,
arXiv:1411.0935 (2014).
\url{https://arxiv.org/abs/1411.0935}.

\bibitem{IC}
IC matroid generator, source repository.
\url{https://github.com/gmou3/matroid-generator}.
Version used: \texttt{60645ae024e831fba3ebd35b2388515b0304140a}.

\bibitem{Ganak}
Ganak exact model counter, version 2.6.4.
\url{https://github.com/meelgroup/ganak}.

\bibitem{Catalogue}
\emph{Matroids with 9 elements}, reference catalogue, Zenodo.
\url{https://zenodo.org/records/6825419}.

\bibitem{Code}
ainta, \emph{Counting all matroids on ten elements}, source repository.
\url{https://github.com/ainta/matroid-count}.
Version used: \texttt{d1aaa1468a01d983b14651a58874948dd357c606}.
\end{thebibliography}
\end{document}
